% CREACION DEL DOCUMENTO
\documentclass[
	spanish,
	letterpaper, oneside
]{article}

% INFORMACION DEL DOCUMENTO
\def\documenttitle {Lista de verificación para evaluar textos digitales}
\def\documentsubtitle {Análisis de contenido de un texto académico digital}
\def\documentsubject {Actividad 4 - Redacción en contextos virtuales}

\def\documentauthor {Martín Jonathan de la Cruz}
\def\coursename {Redacción en contextos virtuales}
\def\coursecode {017}

\def\universityname {Universidad Abierta y a Distancia de México}
\def\universityfaculty {Licenciatura en Derecho}
\def\universitydepartment {Redacción en contextos virtuales}
\def\universitydepartmentimage {departamentos/UnADM}
\def\universitydepartmentimagecfg {height=1.57cm}
\def\universitylocation {Roma Norte, Ciudad de México}

% INTEGRANTES, PROFESORES Y FECHAS
\def\authortable {
	\begin{tabular}{ll}
		\textbf{Alumno:}
		& \begin{tabular}[t]{l}
			 Martín Jonathan de la Cruz \\
		\end{tabular} \\
		Matrícula: & ES2611202040 \\

		\textit{Profesora:}
		& \begin{tabular}[t]{l}
			Ana Rosa \\ Arceo Luna
		\end{tabular} \\
		& \\
		\multicolumn{2}{l}{Fecha de realización: \today} \\
		\multicolumn{2}{l}{\universitylocation}
	\end{tabular}
}

% IMPORTACION DEL TEMPLATE
\input{template}

% FORMATO
\usepackage{helvet}
\renewcommand{\familydefault}{\sfdefault}
\usepackage{array}
\usepackage{longtable}
\usepackage{booktabs}
\usepackage{amssymb}
\setcitestyle{authoryear,open={(},close={)}}

% ==============================================================================
% MARCA DE AGUA EN PORTADA
% ==============================================================================
\def\coverwatermarkenabled {true}
\def\coverwatermarkimage {img/departamentos/UnADM.pdf}
\def\coverwatermarkopacity {0.16}
\def\coverwatermarkwidth {0.72\paperwidth}
\def\coverwatermarkxshift {0cm}
\def\coverwatermarkyshift {-0.35cm}

\newcommand{\insertcoverwatermark}{%
	\ifthenelse{\equal{\coverwatermarkenabled}{true}}{%
		\AddToShipoutPictureBG*{%
			\begin{tikzpicture}[remember picture,overlay]
				\node[
					opacity=\coverwatermarkopacity,
					inner sep=0pt,
					xshift=\coverwatermarkxshift,
					yshift=\coverwatermarkyshift
				] at (current page.center) {%
					\includegraphics[width=\coverwatermarkwidth]{\coverwatermarkimage}%
				};
			\end{tikzpicture}%
		}%
	}{}%
}

\begin{document}

% PORTADA
\insertcoverwatermark
\templatePortrait

% CONFIGURACION DE PAGINA Y ENCABEZADOS
\templatePagecfg
\onehalfspacing

% RESUMEN
\begin{abstractd}
	\textbf{Objetivo:} Evaluar la calidad de un texto académico digital mediante una lista de verificación enfocada en propósito, claridad, estructura, tono, legibilidad y corrección.

	\textbf{Producto elaborado:} Lista de verificación aplicada al texto propio \textit{Comunicación efectiva en el ámbito laboral: juego de roles}, correspondiente a la actividad anterior de la asignatura, más un análisis breve con fortalezas, áreas de mejora y ajustes concretos.

	\textbf{Enfoque de análisis:} La escritura digital no se mide únicamente por estar bien redactada; también debe ser clara, navegable y funcional para quien la lee en pantalla.
\end{abstractd}

% TABLA DE CONTENIDOS
\templateIndex

% CONFIGURACIONES FINALES
\templateFinalcfg

% ======================= INICIO DEL DOCUMENTO =======================

\section{Introducción}

Evaluar un texto digital implica revisar si realmente funciona para el lector. No basta con que el contenido sea correcto; también debe tener propósito claro, estructura ordenada, tono adecuado y una presentación que facilite la lectura en pantalla. En los entornos virtuales, la redacción académica se vuelve una herramienta de organización: guía la interpretación, reduce ambigüedades y permite que el mensaje llegue con mayor precisión.

Desde esta perspectiva, la lista de verificación permite revisar un texto con criterios observables. La escritura académica exige ordenar ideas, seleccionar información pertinente y adecuar el lenguaje al contexto comunicativo \citep{nunezcortesEscrituraAcademicaTeoria2015}. Además, en la universidad, escribir no significa solo cumplir con una entrega, sino construir una explicación comprensible, coherente y sustentada \citep{correaManualBasicoEscritura2018,vitoriaEscrituraAcademicaFormacion2019}. Por ello, en esta actividad se evalúa un texto académico propio de la actividad anterior, con el fin de detectar fortalezas, áreas de mejora y ajustes concretos para hacerlo más claro en formato digital.

\section{Texto digital seleccionado}

El texto seleccionado para la evaluación es un trabajo propio elaborado previamente en la asignatura: \textit{Comunicación efectiva en el ámbito laboral: juego de roles}. Se trata de un documento académico digital de más de una página de extensión, integrado por portada, resumen, introducción, análisis comparativo de correos, juego de roles, reflexión final, conclusión y referencias.

\begin{longtable}{@{}p{0.27\linewidth}p{0.66\linewidth}@{}}
\caption{Ficha del texto digital evaluado}\label{tab:ficha-texto}\\
\toprule
\textbf{Elemento} & \textbf{Descripción} \\
\midrule
\endfirsthead
\toprule
\textbf{Elemento} & \textbf{Descripción} \\
\midrule
\endhead
Tipo de texto & Texto académico digital propio. \\
\addlinespace
Título & \textit{Comunicación efectiva en el ámbito laboral: juego de roles}. \\
\addlinespace
Contexto & Actividad anterior de la materia \textit{Redacción en contextos virtuales}. \\
\addlinespace
Propósito & Analizar correos laborales y reformular un mensaje inadecuado en una versión formal, clara y profesional. \\
\addlinespace
Extensión & Documento mayor a una página, con desarrollo, producto solicitado y cierre reflexivo. \\
\bottomrule
\end{longtable}

\section{Lista de verificación aplicada}

\begin{longtable}{@{}p{0.18\linewidth}p{0.43\linewidth}p{0.13\linewidth}p{0.20\linewidth}@{}}
\caption{Lista de verificación para evaluar el texto digital}\label{tab:lista-verificacion}\\
\toprule
\textbf{Categoría} & \textbf{Criterio} & \textbf{Cumple} & \textbf{Observación breve} \\
\midrule
\endfirsthead
\toprule
\textbf{Categoría} & \textbf{Criterio} & \textbf{Cumple} & \textbf{Observación breve} \\
\midrule
\endhead
Propósito y claridad & El texto tiene un propósito claro. & Sí & El objetivo se identifica desde el resumen y la introducción. \\
\addlinespace
Propósito y claridad & La idea principal se identifica fácilmente desde el inicio. & Sí & La actividad se enfoca en comunicación laboral escrita y correos profesionales. \\
\addlinespace
Organización y estructura & Los párrafos son breves y están bien delimitados. & Parcial & Algunos párrafos concentran varias ideas y pueden dividirse para lectura en pantalla. \\
\addlinespace
Organización y estructura & La información se presenta de forma lógica y ordenada. & Sí & La secuencia avanza de introducción a análisis, producto, reflexión y conclusión. \\
\addlinespace
Organización y estructura & Se utilizan conectores que dan coherencia al texto. & Sí & Hay transiciones como ``desde esta perspectiva'', ``en cambio'' y ``por ello''. \\
\addlinespace
Lenguaje y tono & El lenguaje es adecuado al contexto académico o profesional. & Sí & El vocabulario mantiene formalidad y relación con la práctica laboral. \\
\addlinespace
Lenguaje y tono & Se evita el uso de expresiones coloquiales o informales. & Sí & No se usan expresiones coloquiales dentro del análisis académico; solo aparecen en el correo inadecuado como parte del ejemplo. \\
\addlinespace
Lenguaje y tono & El tono es claro, formal y respetuoso. & Sí & La postura es crítica, pero no agresiva; mantiene intención académica. \\
\addlinespace
Legibilidad en formato digital & El texto es fácil de leer en pantalla. & Parcial & La división por secciones ayuda, aunque algunos bloques podrían hacerse más ligeros. \\
\addlinespace
Legibilidad en formato digital & No hay párrafos excesivamente largos. & Parcial & Existen párrafos amplios que pueden afectar la lectura rápida. \\
\addlinespace
Legibilidad en formato digital & La estructura facilita la lectura rápida o escaneabilidad. & Parcial & Los títulos ayudan, pero se podrían agregar síntesis o viñetas estratégicas. \\
\addlinespace
Corrección & No presenta errores ortográficos. & Sí & No se detectan faltas ortográficas relevantes. \\
\addlinespace
Corrección & No presenta errores gramaticales. & Sí & La sintaxis es clara y coherente. \\
\bottomrule
\end{longtable}

\section{Análisis breve del texto evaluado}

El texto evaluado presenta un propósito claro: analizar la comunicación laboral escrita mediante el contraste entre un correo inadecuado y una versión formal. Su fortaleza principal es la organización por secciones, porque permite ubicar introducción, análisis, juego de roles, reflexión y conclusión sin perder el hilo. También mantiene un tono académico adecuado y explica por qué la claridad, la cortesía y la estructura son condiciones de eficacia comunicativa. Sin embargo, el área de mejora está en la legibilidad digital: algunos párrafos son largos y podrían dificultar la lectura en pantalla. Además, aunque la argumentación es sólida, se podría hacer más escaneable con subtítulos breves o una síntesis visual antes de la reflexión final. Como ajustes concretos, primero dividiría los párrafos extensos en bloques más cortos, cada uno con una función específica: contexto, análisis y consecuencia. Segundo, agregaría una mini-lista de criterios aplicados al correo formal, para que el lector identifique rápido qué cambió y por qué. En general, el texto cumple, pero puede ganar fuerza si reduce densidad y mejora navegación visual.

\section{Ajustes concretos propuestos}

\begin{enumerate}
	\item \textbf{Dividir párrafos extensos.} El texto puede conservar su análisis, pero separando ideas largas en bloques más breves. Esto mejora la legibilidad digital sin perder profundidad.
	\item \textbf{Agregar una síntesis visual.} Antes de la reflexión final, conviene incluir una mini-lista o cuadro breve con tres criterios aplicados: claridad, tono y estructura. Así el lector identifica rápido el cambio entre el correo inadecuado y el formal.
\end{enumerate}

\section{Conclusión}

La lista de verificación muestra que el texto evaluado cumple con su propósito académico y mantiene un tono adecuado para la asignatura. Su estructura general es lógica y permite comprender el análisis de los correos laborales. No obstante, el formato digital exige algo más que buena redacción: también pide navegación visual, párrafos manejables y puntos de apoyo para que el lector pueda ubicar rápido las ideas principales.

Mi conclusión es que el texto funciona, pero todavía puede optimizarse. En una entrega académica digital, la claridad no depende únicamente de tener ideas correctas; depende de cómo se ordenan, se separan y se hacen visibles. La mejora principal no sería cambiar el fondo del trabajo, sino ajustar su presentación para que el análisis sea más directo, más escaneable y más cómodo de leer en pantalla.

\bibliography{RedaccionContextosVirtuales}

\end{document}
